\section{Вращение Земли и геометрия запуска}
\label{sec:launch-geometry}

В пространственной задаче плоскость орбиты согласуют с местом,
направлением и временем старта.

\subsection{Вращение Земли и две скорости}

Ракета на площадке и атмосфера вращаются вместе с Землёй.
Поэтому начальная инерциальная скорость ракеты и её скорость
относительно воздуха без ветра равны
\begin{equation}
\begin{aligned}
 \vect V_0&=\vect\omega_E\times\vect R_0,
 \qquad V_{E,0}=\omega_E R_E\cos\phi,\\
 \vect V_{\mathrm{air}}&=\vect V-\vect\omega_E\times\vect R.
\end{aligned}
\label{eq:launch-initial-velocity}
\end{equation}
Здесь $\phi$ --- геоцентрическая широта. Скорость площадки направлена
на восток и на экваторе составляет около $465$~м/с.
Орбиту определяет $\vect V$, аэродинамику --- $\vect V_{\mathrm{air}}$.
Из векторного сложения видно, что различаются их модули и азимуты.

\courseplot[.55]{geometry-rotation}{Инерциальная и воздушная скорости
без ветра. Скорость атмосферы направлена на восток;
$\vect V=\vect V_{\mathrm{air}}+\vect\omega_E\times\vect R$.}

\subsection{Азимут и наклонение}

Азимут $A$ отсчитываем от севера к востоку: $A=0$ соответствует
движению на север, $A=90^\circ$ --- на восток.
Пусть $\vect e_E$ и $\vect e_N$ --- местные направления востока и севера,
а $\vect k$ --- единичный вектор земной оси.
Единичное горизонтальное направление движения обозначено
на рисунке через $\vect e_A$.
Радиальная скорость не вносит вклада в орбитальный момент,
поэтому нормаль орбиты определяется горизонтальным направлением движения:
\begin{equation}
 \vect n=\frac{\vect R\times\vect V}{\norm{\vect R\times\vect V}}
        =\sin A\,\vect e_N-\cos A\,\vect e_E.
 \label{eq:orbit-normal-azimuth}
\end{equation}
В местном базисе $\vect k=\sin\phi\,\vect e_r+\cos\phi\,\vect e_N$.
Скалярное произведение нормали на земную ось даёт
\begin{equation}
 \boxed{\cos i=\vect n\cdot\vect k=\cos\phi\sin A.}
 \label{eq:launch-azimuth-inclination}
\end{equation}
Так широта и направление движения задают наклонение орбиты.
Здесь $A$ --- \emph{инерциальный} азимут, а $\phi$ относится
к точке, в которой определены положение и скорость.
Подстановка широты космодрома даёт исходную геометрическую оценку
для запуска в плоскость, проходящую через него~\cite{teofilatto}.

\courseplot[.96]{geometry-launch-plane}{Связь местного направления
движения с плоскостью орбиты: проекция её нормали на земную ось
равна $\cos i=\cos\phi\sin A$.}

\subsection{Время старта и плоскость орбиты}

Одно наклонение ещё не фиксирует плоскость: при том же $i$
возможны разные долготы восходящего узла $\Omega$.
Космодром вращается вместе с Землёй, а целевая плоскость
фиксирована в инерциальных осях. Поэтому для прямого запуска
нужно выбрать момент, когда космодром окажется в этой плоскости:
\begin{equation}
 \boxed{\vect n_*\cdot\vect R_s(t_0)=0.}
 \label{eq:launch-window-plane}
\end{equation}
Здесь $\vect n_*$ --- нормаль целевой плоскости,
$\vect R_s(t)$ --- инерциальное положение космодрома.
На рисунке этому условию соответствуют пересечения суточной
параллели с плоскостью орбиты; $\Lambda(t)$ обозначает
инерциальную долготу космодрома.

\courseplot[.92]{geometry-launch-window}{Геометрия выбора времени
старта. В экваториальной проекции космодром движется по окружности;
положения $A$ и $B$ соответствуют его прохождению через целевую плоскость.}

\FloatBarrier
\begin{explanation}{Что задаёт геометрия запуска}
Широта и инерциальный азимут определяют наклонение,
а время старта позволяет выбрать ориентацию плоскости.
Это исходная геометрическая оценка: полная программа выведения
должна обеспечить конечное попадание в плоскость с учётом
вращения Земли и движения ракеты во время полёта.
\end{explanation}
\FloatBarrier
